\documentclass[11pt,letterpaper]{article}

\usepackage[margin=1in]{geometry}
\usepackage[T1]{fontenc}
\usepackage{lmodern}
\usepackage{parskip}
\usepackage{booktabs}
\usepackage{longtable}
\usepackage{array}
\usepackage{enumitem}
\usepackage{hyperref}
\usepackage{xurl}
\usepackage{xcolor}
\usepackage{tabularx}

\hypersetup{
    colorlinks=true,
    linkcolor=blue!60!black,
    urlcolor=blue!60!black
}

\setlength{\emergencystretch}{2em}

\newcommand{\notready}{\textbf{\textcolor{red!75!black}{NOT READY}}}
\newcommand{\ready}{\textbf{\textcolor{green!50!black}{READY}}}

\title{\textbf{Ruthless Technical Audit of the Polymarket Insider-Signal SRS}}
\author{Architecture and Machine Learning Readiness Review}
\date{April 7, 2026}

\begin{document}
\maketitle
\tableofcontents
\newpage

\section{Executive Verdict}

\textbf{Binary verdict: \notready.}

This SRS is a strong research memo and a weak production build specification. It contains a credible research question, several useful implementation directions, and a correct instinct that wallet-aware flow matters. It is \emph{not} implementation-ready because it still hand-waves the three hardest parts of the system:

\begin{enumerate}[nosep]
    \item \textbf{data semantics} (what exactly is observable, replayable, and point-in-time valid),
    \item \textbf{stream-state architecture} (how online state, dedupe, late events, and replay actually work), and
    \item \textbf{operational economics} (API fan-out, storage growth, and always-on runtime cost).
\end{enumerate}

If you started coding directly from this SRS, the plan would fail in production for predictable reasons: incomplete identity semantics, non-replayable event handling, fragile serverless streaming assumptions, and a cost model that is materially understated.

\textbf{Repository evidence that the plan is not close to model-ready:}
\begin{itemize}[nosep]
    \item \texttt{run\_collector.py} still hard-filters the market universe to sports tags.
    \item \texttt{collectors/trades\_collector.py} still calls the trades API with a token-driven parameter path instead of the condition-level semantics described in the SRS.
    \item \texttt{database/schema.sql} still stores trades without \texttt{proxy\_wallet}, \texttt{transaction\_hash}, \texttt{condition\_id}, or \texttt{outcome\_side}.
    \item \texttt{collectors/ttl\_manager.py} still deletes closed-market trades and snapshots after 24 hours.
    \item \texttt{ml\_pipeline/feature\_builder.py} is still a stub.
\end{itemize}

The most damaging strategic mistake in the document is this: it treats data collection, online detection, external evidence retrieval, and backtesting as one stack. They are not one stack. They are four systems with different latency, durability, and correctness requirements.

\section{1. Reality Check: Assumptions vs.\ Concrete Steps}

\subsection{Material Assumptions Extracted}

The table below lists the material assumptions that affect buildability. These are the assumptions that will break the system, not merely stylistic or wording assumptions.

\begingroup
\footnotesize
\setlength{\tabcolsep}{4pt}
\begin{longtable}{p{0.7cm}p{3.7cm}p{5.0cm}p{5.0cm}}
\toprule
\textbf{ID} & \textbf{Assumption} & \textbf{Why This Is Fragile} & \textbf{Critical Missing Step} \\
\midrule
\endfirsthead
\toprule
\textbf{ID} & \textbf{Assumption} & \textbf{Why This Is Fragile} & \textbf{Critical Missing Step} \\
\midrule
\endhead
A1 & Public Polymarket feeds provide enough information to estimate market microstructure toxicity, quote replenishment, and informed flow. & Public market data exposes price, trades, and book state, but not the full order lifecycle. You do not directly observe every cancel, repost, or queue-position event from the public side. & Separate ``observable from public feed'' features from ``unobservable without order lifecycle'' features and explicitly drop anything you cannot estimate honestly. \\
\addlinespace
A2 & A \texttt{proxyWallet} is a useful identity primitive. & Proxy wallets are not people. One human can use multiple wallets, and one wallet can represent delegated or shared execution patterns. & Define an entity-resolution layer and a strict policy that the system scores wallets or clusters, not humans. \\
\addlinespace
A3 & ``First seen in local history'' is a valid definition of a new wallet. & Collector downtime, stale history, partial market coverage, and TTL deletion all manufacture fake novelty. ``New to your database'' is not the same as ``new to the platform.'' & Formalize first-seen semantics over an immutable full archive and attach provenance to every first-seen claim. \\
\addlinespace
A4 & Condition-level enrichment via \texttt{/trades}, \texttt{/activity}, \texttt{/positions}, and profile endpoints can run online at scale. & The fan-out grows with markets, wallets, and event bursts. Full-universe enrichment will hit quotas, latency walls, and queue backlog. & Define enrichment budgets, caching policy, prioritization rules, and a deferred enrichment queue. \\
\addlinespace
A5 & REST and WebSocket timestamps are directly comparable. & Arrival time, source timestamps, retries, and API lag create out-of-order events. This destroys naive episode reconstruction. & Define event time, processing time, watermarks, clock-skew handling, and dedupe keys. \\
\addlinespace
A6 & ``Public silence'' is measurable from X plus web/news search. & Missing retrieval coverage, indexing lag, provider outages, and bad queries all look like silence. The score can end up measuring search blindness, not public ignorance. & Build a retrieval benchmark with timestamp normalization, source coverage accounting, and point-in-time caches. \\
\addlinespace
A7 & Cloud Run is an easy production host for a 24/7 stateful WebSocket collector. & Cloud Run WebSocket connections are still request-bounded, session affinity is only best effort, and open sockets are billed as active instance time.\footnote{\url{https://cloud.google.com/run/docs/triggering/websockets}}\footnote{\url{https://cloud.google.com/run/docs/configuring/request-timeout}}\footnote{\url{https://cloud.google.com/run/pricing}} & Design reconnect choreography, partition ownership, state externalization, and a real cost model before picking the runtime. \\
\addlinespace
A8 & PostgreSQL plus Redis is sufficient even without an append-only event log. & Neither database is a substitute for a durable replayable event stream. Without a log, failures become silent data holes and backtests become fiction. & Add a durable event log and explicit replay pipeline before any serious detector is built. \\
\addlinespace
A9 & Raw payload archival can be added later without hurting research quality. & If raw events are not archived from day one, you cannot reconstruct feature logic, drift, or point-in-time state later. & Define an immutable raw archive format now, including schema versioning and partitioning. \\
\addlinespace
A10 & The listed labels are operationally well-defined. & ``Next-5m return,'' ``public corroboration lag,'' and ``event outcome'' are different tasks with different censoring and leakage risks. & Write a formal label specification with horizons, censoring rules, boundary behavior, and market-state exclusions. \\
\addlinespace
A11 & The proposed metrics and thresholds are meaningful. & AUROC, Brier, precision@k, paper-trade PnL, and Sharpe-like ratios answer different questions. The SRS mixes ranking, classification, and trading objectives without priority. & Establish a metric hierarchy: first ranking utility, then calibration, then economic lift, then operational noise. \\
\addlinespace
A12 & A simple best-bid/best-ask simulator is enough for paper trading. & Queue position, partial fills, fees, slippage curvature, and adverse selection dominate short-horizon execution. & Build a conservative execution simulator that models fill probability, fees, partial fills, and liquidity caps. \\
\addlinespace
A13 & Graph models and sequence models belong early in the roadmap. & Complex models amplify garbage when identity semantics and replayability are not stable. Right now the data contract is not mature enough. & Delay complex sequence/graph models until the feature store, label taxonomy, and entity semantics are stable. \\
\addlinespace
A14 & The infrastructure cost is approximately \$40--75/month. & This is materially understated. Open WebSocket Cloud Run services are billed as active compute, a 1 GiB Basic Redis instance alone is already far above the document's \$0--10 estimate, and the smallest Cloud SQL shared-core instances are explicitly not production-oriented.\footnote{\url{https://cloud.google.com/memorystore/pricing}}\footnote{\url{https://cloud.google.com/sql/docs/postgres/instance-settings}} & Rebuild the cost model to include active compute, Redis floor cost, storage growth, logging, network egress, API spend, and backfill jobs. \\
\addlinespace
A15 & The development roadmap is realistic. & The roadmap treats ingestion correctness, stateful streaming, external evidence retrieval, backtesting, and thesis-grade validation as a linear sequence with optimistic durations. They are overlapping projects with dependencies. & Replace the roadmap with hard stage gates and explicit exit criteria. \\
\addlinespace
A16 & The Google Custom Search dependency is correctly understood. & The SRS appears to conflate different Google search APIs. Google's current documentation still describes the Programmable Search JSON API; the deprecation note in the SRS is not a safe foundation for architecture decisions.\footnote{\url{https://developers.google.com/custom-search/v1/overview}} & Make vendor dependency claims only from current vendor docs and record the exact API being referenced. \\
\addlinespace
A17 & Current repository state is close enough to justify immediate modeling work. & It is not. The collector still has sports-only filtering, one-sided live coverage, missing wallet fields, TTL deletion, and no feature builder. & Finish the data-correctness phase before writing nontrivial ML code. \\
\addlinespace
A18 & Legal/compliance risk is mostly a wording problem. & It is also an access-control, audit-log, retention, and operational-governance problem. & Add security, secret-management, access-control, and auditability requirements. \\
\bottomrule
\end{longtable}
\endgroup

\subsection{What Is Actually Concrete}

The SRS is not devoid of implementation detail. The following items are concrete enough to start engineering from:

\begin{itemize}[nosep]
    \item Remove the sports-only universe filter or make it configurable.
    \item Fix trade ingestion so the collector uses condition-level identifiers correctly.
    \item Subscribe to both YES and NO assets instead of only one side.
    \item Add wallet-linked trade fields: \texttt{proxy\_wallet}, \texttt{transaction\_hash}, \texttt{condition\_id}, \texttt{asset}, \texttt{outcome\_side}, \texttt{price}, \texttt{size}, \texttt{usdc\_notional}, \texttt{timestamp}, and \texttt{source}.
    \item Archive raw payloads before normalization.
    \item Preserve closed-market microstructure long enough to support event studies and label construction.
    \item Add wallet/profile/activity/position tables and alert outcome tables.
    \item Separate candidate generation, wallet aggregation, anomaly scoring, evidence retrieval, and alert ranking into explicit stages.
    \item Keep external evidence retrieval off the hot path until a candidate episode exists.
    \item Use point-in-time replay and holdout-based validation instead of random train/test splits.
    \item Use Telegram as primary notification and keep alert history in the database.
\end{itemize}

Those are real steps. They are just not sufficient. They solve feature checklist problems, not system correctness.

\subsection{Where Assumptions Are Masking Missing Engineering Work}

\begin{enumerate}[leftmargin=1.5em]
    \item \textbf{``Fresh-wallet episode'' masks an identity problem.} You cannot operationalize this until you define what a wallet entity is, how you cluster proxies, how you treat re-appearing wallets after downtime, and how you separate ``new to us'' from ``new to market.'' Right now the phrase is statistically seductive and operationally underspecified.
    \item \textbf{``Public-silence score'' masks a retrieval-quality problem.} The missing step is a benchmarked retrieval stack with cached point-in-time results and timestamp provenance. Without that, silence is just missing data with a more flattering name.
    \item \textbf{``Point-in-time reconstruction'' masks an event-log problem.} You need immutable raw archives, event-time watermarks, model-version tracking, feature-version tracking, and replay tooling. None of that is specified.
    \item \textbf{``Cloud migration'' masks a runtime-selection problem.} The document names services, but it does not define partition ownership, failover behavior, dedupe semantics, or where streaming state lives during reconnects.
    \item \textbf{``Alert ranking'' masks a labeling problem.} Ranking requires analyst feedback loops, clear rejection reasons, and calibrated false-positive accounting. The SRS names \texttt{accepted/ignored/rejected} labels but does not specify the workflow that produces them.
    \item \textbf{``Paper trading'' masks a market-impact problem.} The document does not specify fill simulation, fee treatment, or exposure caps. Without those, paper-trade performance will be fantasy.
    \item \textbf{``Goodhart validation'' masks an observability instrumentation problem.} You cannot test observability shocks unless you explicitly record dashboard publication, alert publication, media exposure, and audience reach.
    \item \textbf{``Sequence/graph models later'' masks a data-contract maturity problem.} These models require stable semantics. Today the data contract is still missing the fields that would make the graphs and sequences meaningful.
\end{enumerate}

\section{2. Implementation Feasibility}

\subsection{Architecture Critique}

\textbf{Is the proposed system buildable with current technology?} Yes, but \emph{not} in the exact form implied by the SRS. It is buildable only if you separate the architecture into four planes:

\begin{enumerate}[nosep]
    \item \textbf{Sensing plane} --- durable raw market ingestion,
    \item \textbf{State plane} --- online rolling state and entity state,
    \item \textbf{Decision plane} --- candidate scoring and ranking,
    \item \textbf{Research plane} --- replay, validation, and training.
\end{enumerate}

The SRS currently compresses those planes into a list of services, but a list of services is not an architecture. What is missing is the contract between them.

\textbf{What is structurally wrong with the proposed workflow:}

\begin{itemize}[nosep]
    \item \textbf{PostgreSQL is being treated like a stream processor.} It is an operational store, not your replay log and not your low-latency feature state engine.
    \item \textbf{The evidence engine is underspecified as a real-time component.} X/web retrieval is seconds-scale and quota-limited. It must be asynchronous and best-effort, not part of a sub-second scoring path.
    \item \textbf{The data model names entities but does not define their keys or update semantics.} There is no specification for uniqueness, late-arriving updates, dedupe between REST and WebSocket sources, or replay idempotency.
    \item \textbf{The SRS assumes a hot-path detector can be built before state semantics are formalized.} That is backwards. You need trustworthy state before you need a model.
    \item \textbf{The deployment recommendation is too serverless by default.} Cloud Run can host stateless detectors and batch jobs. It is a compromised choice for a persistent streaming collector because timeouts, reconnects, and best-effort affinity force you to externalize more state than the SRS accounts for.\footnote{\url{https://cloud.google.com/run/docs/triggering/websockets}}
\end{itemize}

\textbf{Corrected architecture that is actually buildable:}

\begin{enumerate}[nosep]
    \item \textbf{Collector worker} writes every raw WebSocket and REST event to immutable object storage and publishes normalized envelopes to a durable queue.
    \item \textbf{Stream state worker} maintains rolling market, wallet, and event-family state in Redis or another explicit online state store.
    \item \textbf{Candidate detector} runs deterministic rules plus lightweight anomaly scores on streaming features.
    \item \textbf{Ranker} loads a calibrated gradient-boosted model in-process and scores only candidate episodes.
    \item \textbf{Evidence worker} asynchronously enriches candidates with X/web/news evidence and writes a versioned evidence snapshot.
    \item \textbf{Alert worker} applies suppression, severity policy, and human-readable rendering.
    \item \textbf{Replay/training jobs} rebuild features from raw archives, never from mutable live tables alone.
\end{enumerate}

\textbf{Brutal truth:} the collector is still in the ``make the data trustworthy'' phase. The SRS is written as if the project is already in the ``rank and alert'' phase.

\subsection{Primary Bottlenecks and Failure Points}

\begin{enumerate}[leftmargin=1.5em]
    \item \textbf{API fan-out and rate limits.} Wallet activity, positions, profiles, and external evidence calls cannot be naively executed for every active market. Selective enrichment is mandatory.
    \item \textbf{State management hot spots.} ``Fresh wallet count in last 15 minutes'' and ``cluster concentration'' require online keyed state with expiry, not ad hoc SQL queries over hot tables.
    \item \textbf{Out-of-order and duplicate events.} The same economic event can appear via WebSocket, REST trade pulls, and later backfills. Without source-priority rules and event ids, you will double count or mis-time signals.
    \item \textbf{Storage write amplification.} You are already sitting on a large local SQLite database while still deleting closed-market data aggressively. Full-universe collection plus raw archives plus wallet enrichments will multiply storage, not add to it incrementally.
    \item \textbf{Backtest leakage.} Wallet profiles, traded counts, and public evidence are slowly changing and time-sensitive. If you do not version them point-in-time, the backtest will leak future information.
    \item \textbf{External evidence latency.} Public search is not millisecond infrastructure. Treating it like a real-time feature source will cause either timeout pressure or silent skipping.
    \item \textbf{Alert fatigue.} The SRS names severity levels but does not define dedupe, cooldowns, escalation policy, or analyst acknowledgment. That produces spam, not monitoring.
    \item \textbf{Cost drift.} The document's operating-cost estimate does not survive contact with always-on active compute, Redis floor cost, storage growth, log volume, and paid external APIs.
\end{enumerate}

\subsection{Implementation Verdict}

\textbf{Verdict: \notready.}

\textbf{Justification:}

\begin{itemize}[nosep]
    \item The system boundary between ingestion, online state, and offline research is not specified rigorously enough to build without rework.
    \item The current repository state is missing the wallet-linked data fields that the core hypotheses depend on.
    \item The document overstates what is measurable from public market data and understates the engineering needed for replayability and latency control.
    \item The runtime and cost assumptions are materially optimistic.
    \item The ML section jumps to models before the data contract, label semantics, and evaluation contract are stable.
\end{itemize}

This is not a ``small polish missing'' situation. This is a ``requirements document still one major layer too high-level'' situation.

\section{3. Machine Learning Strategy and Utilization}

\subsection{What Models You Should Actually Use}

\textbf{Do not start with transformers, graph neural networks, or generic one-class anomaly detectors in production.}

They are tempting because they sound sophisticated. They are the wrong first move here because your real bottleneck is semantic correctness, not expressive model capacity.

\textbf{Recommended model stack by stage:}

\begin{enumerate}[leftmargin=1.5em]
    \item \textbf{Online candidate generation: deterministic rules + robust streaming statistics.}

    Use rolling median/MAD, exponentially weighted z-scores, CUSUM-style changepoint logic, and extreme-value tail thresholds over episode features such as:
    \begin{itemize}[nosep]
        \item fresh-wallet notional share,
        \item fresh-wallet count,
        \item directional imbalance,
        \item probability velocity and acceleration,
        \item depth collapse and spread widening,
        \item concentration ratio and event-family spread.
    \end{itemize}

    \textbf{Reason:} this gives you online recall with interpretable failure modes. It is also robust to label scarcity.

    \item \textbf{Primary alert ranker: LightGBM or CatBoost over episode-level features.}

    The first serious learned model should be a tabular ranker/classifier over candidate episodes, not raw ticks. Use separate horizons or multi-head outputs for 5m, 30m, and 2h forward movement; then calibrate per liquidity bucket and market family.

    \textbf{Use:}
    \begin{itemize}[nosep]
        \item LightGBM if you want maximum training speed and straightforward feature importance,
        \item CatBoost if categorical handling and missing-value behavior become dominant.
    \end{itemize}

    \textbf{Do not use:} raw deep transformers as the first production ranker.

    \item \textbf{Temporal dynamics model: marked Hawkes process first, TCN second.}

    This domain is bursty, event-driven, and irregularly timed. A marked Hawkes process is a better conceptual fit for trade-arrival intensity and coordinated bursts than a generic transformer. Once the data contract is stable, a temporal convolutional network (TCN) over 1s or 5s aggregated sequences is the sensible next step.

    \textbf{Do not start with:} generic sequence transformers. They are expensive, hard to calibrate, and unnecessary at this maturity level.

    \item \textbf{Graph layer: graph features and community detection before any GNN.}

    Build temporal wallet-market graphs and compute:
    \begin{itemize}[nosep]
        \item degree, weighted degree, and concentration features,
        \item event-family overlap,
        \item co-trade synchronization,
        \item repeated same-direction participation,
        \item cluster persistence across markets.
    \end{itemize}

    Use Leiden or Louvain-style community detection, or HDBSCAN over graph-derived embeddings or pairwise features, before even thinking about GraphSAGE or TGN.

    \item \textbf{External evidence layer: retrieval and reranking, not LLM-first scoring.}

    Use a classic retrieval stack:
    \begin{itemize}[nosep]
        \item lexical retrieval (BM25 or equivalent),
        \item dense retrieval or embeddings for semantic expansion,
        \item cross-encoder or reranker for top documents,
        \item optional small LLM only for summary generation after ranking.
    \end{itemize}

    The LLM should explain a candidate, not decide whether the candidate exists.
\end{enumerate}

\subsection{How the ML Stack Should Integrate With the Backend}

\textbf{Critical design rule: no heavyweight ML model belongs on the raw-event hot path.}

The integration should look like this:

\begin{enumerate}[nosep]
    \item \textbf{Ingestion worker} receives raw market events and writes immutable raw records.
    \item \textbf{Feature-state worker} updates rolling features in memory/Redis keyed by market, event family, wallet, and wallet-cluster.
    \item \textbf{Candidate detector} emits an episode only when deterministic and anomaly thresholds pass.
    \item \textbf{Episode ranker} loads the LightGBM/CatBoost model in-process and scores the episode in milliseconds on CPU.
    \item \textbf{Evidence worker} performs X/web/news retrieval asynchronously and enriches the candidate. It does not block the initial WATCH-level alert if the evidence fetch is slow.
    \item \textbf{Alert composer} renders the final packet with evidence, invalidation conditions, and confidence calibration.
\end{enumerate}

\textbf{Recommended latency budget:}

\begin{center}
\begin{tabular}{@{}p{4.8cm}p{3.2cm}p{6.0cm}@{}}
\toprule
\textbf{Stage} & \textbf{Budget} & \textbf{Notes} \\
\midrule
Raw event ingest to durable log & $< 250$ ms median & This is a systems problem, not an ML problem. \\
Feature-state update & $< 1$ s median & Online rolling windows and keyed state only. \\
Candidate rule + anomaly score & $< 50$ ms & CPU only. \\
Episode ranker inference & $< 10$ ms & LightGBM/CatBoost in-process. \\
Evidence retrieval & 2--10 s best-effort & Async, timeout-bounded, not on the critical path for candidate creation. \\
Alert dispatch & $< 1$ s & Telegram/Discord webhook class latency. \\
\bottomrule
\end{tabular}
\end{center}

\textbf{Serving recommendation:}
\begin{itemize}[nosep]
    \item Keep the ranker inside the detector process until throughput proves otherwise.
    \item Retrain offline on replayed episodes, not live mutable tables.
    \item Version every model artifact, feature schema, and threshold set.
    \item Store every emitted score with the feature vector hash or feature-schema version needed to reproduce it later.
\end{itemize}

\subsection{Hidden Data, Compute, and Cost Requirements the SRS Ignores}

\begin{enumerate}[leftmargin=1.5em]
    \item \textbf{Streaming state is the main compute cost, not the ranker.} A LightGBM model is cheap. Maintaining rolling wallet and market state with correctness is not.
    \item \textbf{Storage growth is already warning you.} The local SQLite database is already large while coverage is incomplete and TTL deletion is aggressive. Full-universe, wallet-aware, raw-archived collection will produce materially higher storage and I/O costs.
    \item \textbf{Cloud Run active WebSocket cost is not idle-cost math.} Open WebSocket connections are billed as active time, not as a sleepy cron job.\footnote{\url{https://cloud.google.com/run/docs/triggering/websockets}}\footnote{\url{https://cloud.google.com/run/pricing}}
    \item \textbf{Redis is not a rounding error.} The SRS estimate of \$0--10/month for Redis-class infrastructure is not credible given current Memorystore pricing.\footnote{\url{https://cloud.google.com/memorystore/pricing}}
    \item \textbf{Shared-core Cloud SQL is not thesis-grade ingestion infrastructure.} Google's own docs position the smallest shared-core instances as non-production-oriented.\footnote{\url{https://cloud.google.com/sql/docs/postgres/instance-settings}}
    \item \textbf{External evidence costs are missing.} X API spend, web/news search provider spend, and query retry volume are not included.
    \item \textbf{Backtesting is a data-engineering workload.} Replaying months of raw events and recomputing features point-in-time will cost CPU, storage, and operational attention, even if model training itself stays modest.
    \item \textbf{If you add LLM summarization or dense embedding generation, you add either API cost or GPU cost.} That is fine if intentional; the SRS currently treats it like free glue logic.
\end{enumerate}

\textbf{Bottom line on ML:} this project should start as a streaming statistics and gradient-boosted ranking system. If the data becomes trustworthy, then sequence and graph models become legitimate. If the data remains semantically weak, deep models will only give you more expensive hallucinations.

\section{4. Actionable Improvements Before Writing More Code}

\subsection{Priority-Ordered Fixes}

\begin{itemize}[nosep]
    \item \textbf{P0: Formalize the architecture around a durable event log.} Do not keep pretending PostgreSQL is enough. Add an append-only raw event stream and a replay contract.
    \item \textbf{P0: Finish ingestion correctness before ML work.} Remove the sports filter, fix condition-level trade ingestion, capture both assets, stop deleting research-critical closed-market data before it is archived, and add wallet-linked trade fields.
    \item \textbf{P0: Define time semantics.} Specify event time, processing time, watermark policy, dedupe keys, and late-event handling.
    \item \textbf{P0: Define entity semantics.} Specify what a wallet, wallet profile, cluster, and first-seen event actually mean, and how they are versioned over time.
    \item \textbf{P0: Replace the serverless hand-wave with a runtime decision memo.} Decide whether the persistent collector lives on a VM, container worker, or another long-lived runtime, and justify it with timeout, state, and cost math.
    \item \textbf{P0: Rewrite the cost model.} Include active Cloud Run time, Redis floor cost, Cloud SQL sizing, storage growth, logging, egress, search APIs, and backfill runs.
    \item \textbf{P1: Write a label specification.} Define targets, horizons, censoring, exclusions near resolution, and the exact mapping from score to operational action.
    \item \textbf{P1: Build the external evidence benchmark before productizing ``public silence.''} Measure retrieval recall, provider lag, timestamp reliability, and cache completeness.
    \item \textbf{P1: Define the alert feedback loop.} Specify analyst review, rejection reasons, suppression windows, dedupe policy, and how feedback becomes labels.
    \item \textbf{P1: Build a conservative execution simulator.} Include fees, queue position assumptions, liquidity caps, and adverse selection.
    \item \textbf{P2: Delay graph and deep temporal models until the feature store is stable.} Use tabular ranking and graph features first.
    \item \textbf{P2: Add governance requirements.} Secrets, access control, audit logging, and retention policy belong in the SRS, not as afterthoughts.
\end{itemize}

\subsection{The Minimum Hard Gate for ``Implementation Ready''}

This document should not be called implementation-ready until all of the following are true:

\begin{enumerate}[nosep]
    \item raw event archival is designed and implemented,
    \item online state semantics are defined,
    \item wallet-linked ingestion is correct and durable,
    \item label definitions are written,
    \item the runtime and cost model are numerically credible, and
    \item at least one end-to-end replay from raw data to alert can be reproduced deterministically.
\end{enumerate}

Until then, the correct status is not ``start building the full system.'' The correct status is ``finish the data and state contract so the system can be built once instead of twice.''

\section*{External Sources Checked}
\addcontentsline{toc}{section}{External Sources Checked}

\begin{itemize}[nosep]
    \item Polymarket WebSocket overview: \url{https://docs.polymarket.com/quickstart/websocket/wss-overview}
    \item Polymarket market channel quickstart: \url{https://docs.polymarket.com/quickstart/websocket/market-channel}
    \item Polymarket trades API reference: \url{https://docs.polymarket.com/api-reference/core/get-trades-for-a-user-or-markets}
    \item Cloud Run WebSockets: \url{https://cloud.google.com/run/docs/triggering/websockets}
    \item Cloud Run request timeout: \url{https://cloud.google.com/run/docs/configuring/request-timeout}
    \item Cloud Run pricing: \url{https://cloud.google.com/run/pricing}
    \item Cloud SQL instance settings: \url{https://cloud.google.com/sql/docs/postgres/instance-settings}
    \item Memorystore pricing: \url{https://cloud.google.com/memorystore/pricing}
    \item Google Programmable Search JSON API overview: \url{https://developers.google.com/custom-search/v1/overview}
    \item X recent search overview: \url{https://docs.x.com/x-api/posts/search/recent-posts/introduction}
\end{itemize}

\end{document}
